% !TEX TS-program = xelatex
% !TEX encoding = UTF-8 Unicode
% -*- coding: UTF-8; -*-
% vim: set fenc=utf-8

\documentclass[11pt,letterpaper]{article}

\usepackage[utf8]{inputenc}
\usepackage[T1]{fontenc}
\usepackage[margin=0.5in]{geometry}
\usepackage{hyperref}
\usepackage{xcolor}
\usepackage{enumitem}

% Define colors
\definecolor{darkblue}{rgb}{0.0, 0.0, 0.5}
\definecolor{darkgray}{rgb}{0.3, 0.3, 0.3}

% Remove page number
\pagestyle{empty}

% Define custom commands
\newcommand{\name}[1]{
  \centerline{\Large\textbf{#1}}
  \vspace{0.5em}
}

\newcommand{\contact}[4]{
  \centerline{\small #1 \quad $\bullet$ \quad #2 \quad $\bullet$ \quad #3 \quad $\bullet$ \quad #4}
  \vspace{1em}
}

\newcommand{\sectiontitle}[1]{
  \vspace{0.3em}
  {\large\textbf{\textcolor{darkblue}{#1}}}
  \vspace{0.2em}
  \hrule
  \vspace{0.3em}
}

\newcommand{\subsectiontitle}[2]{
  \textbf{#1} \hfill \textit{#2}
  \vspace{0.2em}
}

\newcommand{\entry}[4]{
  \textbf{#1} \hfill #2 \\
  \textit{#3} \\
  #4
  \vspace{0.3em}
}

\begin{document}

% Title Block
\name{Le Duc Nguyen}

\contact
  {\href{mailto:nguyen.ld2310@gmail.com}{nguyen.ld2310@gmail.com}}
  {+84 (945) 032-669}
  {\href{https://github.com/AndrewHaward2310}{GitHub}}
  {Hanoi, Vietnam}

% Research Interests
\sectiontitle{RESEARCH INTERESTS}

Intelligent manufacturing systems at the intersection of \textbf{multi-agent reinforcement learning}, \textbf{industrial robotics}, and \textbf{IoT-driven automation}. Focused on developing AI frameworks that bridge cutting-edge research and real-world factory deployment --- enabling human-AI collaborative systems where technology amplifies human capability rather than replacing it. Aspiring to address the digital transformation gap in Southeast Asian manufacturing through inclusive, incremental Industry 4.0 solutions.

% Education
\sectiontitle{EDUCATION}

\entry
  {Hanoi University of Science and Technology (HUST)}
  {Oct 2021 -- Mar 2026}
  {B.S. in Mechatronics Engineering}
  {
    \begin{itemize}[leftmargin=1.5em, topsep=0pt, itemsep=0.2em]
      \item Interdisciplinary training spanning robotics, automation, embedded control, and intelligent systems
      \item Relevant coursework: Automatic Control, Robotics, Digital Systems Design, Signal Processing, Industrial Automation
    \end{itemize}
  }

% Research and Project Experience
\sectiontitle{RESEARCH AND PROJECT EXPERIENCE}

\entry
  {\href{https://andrewhaward2310.github.io/gnn-il-ppo-paper/}{Multi-Agent Reinforcement Learning for AGV Coordination in Smart Factories}}
  {Jul 2025 -- Nov 2025}
  {Project Lead \& Primary Researcher | School of Mechanical Engineering, HUST}
  {
    \begin{itemize}[leftmargin=1.5em, topsep=0pt, itemsep=0.2em]
      \item Proposed a novel multi-agent reinforcement learning (MARL) framework for autonomous guided vehicle (AGV) coordination, combining \textbf{Proximal Policy Optimization (PPO)} with \textbf{Graph Neural Networks (GNN)} over a heterogeneous super-graph representation
      \item Demonstrated that integrating classical PID controllers with modern RL agents yields complementary performance --- validating an \textbf{incremental transition pathway} from legacy control to intelligent systems without requiring full infrastructure replacement
      \item Designed cooperative learning policies enabling multiple AGVs to simultaneously learn task allocation, collision avoidance, and deadlock resolution through training rather than explicit rule-based programming
      \item Developed a high-fidelity 2D simulation environment for benchmarking multi-agent coordination strategies under varying task densities
      \item Team size: 2 | Technologies: PPO, GNN, Imitation Learning, PID Control, ROS, Python
    \end{itemize}
  }

\entry
  {LLM-based Intelligent Fault Query System for Industrial Maintenance}
  {Nov 2023 -- Jan 2024}
  {Project Lead | DENSO Factory Hacks 2023 (FPT $\times$ DENSO Collaboration)}
  {
    \begin{itemize}[leftmargin=1.5em, topsep=0pt, itemsep=0.2em]
      \item Led a cross-organizational team to develop an \textbf{AI-powered fault diagnosis system} that enabled maintenance engineers to query historical error databases using natural language --- reducing manual lookup time from hours to seconds
      \item Designed a hybrid retrieval architecture integrating \textbf{Elasticsearch} for keyword-based indexing and \textbf{Weaviate} (vector database) for semantic search, achieving high-precision answers across large-scale industrial fault records
      \item Gained first-hand exposure to factory-floor operations at DENSO, observing both the enormous potential and the practical barriers of deploying AI in real manufacturing environments
      \item \textbf{Award:} Third Prize, DENSO Factory Hacks 2023
      \item Team size: 4 | Technologies: Large Language Models, Elasticsearch, Weaviate, Docker
    \end{itemize}
  }

\entry
  {Real-time Asset Tracking System using Ultra-Wideband (UWB) Technology}
  {Nov 2022 -- Mar 2023}
  {Developer \& Data Engineer | Innovation and Scientific Research Project}
  {
    \begin{itemize}[leftmargin=1.5em, topsep=0pt, itemsep=0.2em]
      \item Developed a warehouse management system leveraging \textbf{UWB indoor positioning} for real-time asset tracking and location-aware inventory management
      \item Engineered end-to-end data pipeline: UWB sensor data acquisition via MQTT $\rightarrow$ MongoDB storage $\rightarrow$ analytics with Pandas --- establishing foundational experience in IoT-to-cloud data architectures for industrial environments
      \item \textbf{Award:} HUST Innovation and Scientific Research Award, 2022
      \item Team size: 4 | Technologies: UWB Sensors, MQTT, MongoDB, FastAPI, Pandas
    \end{itemize}
  }

\entry
  {\href{https://github.com/AndrewHaward2310/IoTSmartGardenSystem}{IoT-based Intelligent Monitoring and Control System}}
  {Oct 2025 -- Jan 2026}
  {System Architect \& Backend Lead | University Course Project}
  {
    \begin{itemize}[leftmargin=1.5em, topsep=0pt, itemsep=0.2em]
      \item Architected a scalable IoT platform integrating \textbf{edge devices (ESP32)} with cloud data infrastructure for real-time environmental monitoring and automated control --- demonstrating industrial IoT deployment patterns applicable to smart factory scenarios
      \item Designed a multi-protocol data ingestion pipeline (MQTT for sensor streams, REST for management APIs) with a polyglot persistence layer (PostgreSQL, MongoDB, Redis), reflecting real-world industrial data architecture challenges
      \item Team size: 4 | Technologies: Node.js, PostgreSQL, MongoDB, MQTT, Docker, Microservices
    \end{itemize}
  }

% Awards and Scholarships
\sectiontitle{AWARDS AND HONORS}

\subsectiontitle{Third Prize --- DENSO Factory Hacks 2023}{Jan 2024}
Organized by DENSO and FPT | Recognized for innovative LLM-based industrial fault diagnosis system
\vspace{0.3em}

\subsectiontitle{Innovation and Scientific Research Award}{Jun 2022}
Hanoi University of Science and Technology | UWB-based warehouse tracking system
\vspace{0.3em}

\subsectiontitle{Dean's List / Academic Excellence Scholarship}{2021 -- 2022}
School of Mechanical Engineering, HUST | Awarded for high academic performance
\vspace{0.3em}

% Skills
\sectiontitle{TECHNICAL SKILLS}

\textbf{AI \& Machine Learning:} Reinforcement Learning (PPO, MARL, Imitation Learning), Graph Neural Networks (GNN), Large Language Models, Semantic Search (Weaviate), Natural Language Processing

\vspace{0.2em}
\textbf{Robotics \& Control Systems:} ROS (Robot Operating System), PID Control, Motion Planning, Multi-Agent Coordination, Autonomous Guided Vehicles (AGVs), Industrial Automation

\vspace{0.2em}
\textbf{Data \& IoT Systems:} Python (Advanced), MATLAB (Advanced), C/C++, MongoDB, PostgreSQL, MQTT Protocol, Docker, Data Pipeline Design, Real-time Sensor Data Processing

% References
\sectiontitle{REFERENCES}

\textbf{Assoc. Prof. Dr. Pham Duc An}
\newline
Senior Lecturer, School of Mechanical Engineering
\newline
Vice Head, Department of Mechatronics Engineering
\newline
Hanoi University of Science and Technology (HUST)
\newline
\href{mailto:an.phamduc@hust.edu.vn}{an.phamduc@hust.edu.vn}

\end{document}
